\section{Introduction}
Voice is one of the most frequent and naturally used communication interfaces for humans. With recent developments in technology, voice is being used as a main interface in artificial intelligence (AI) voice assistant services such as Amazon Alexa, and it is also widely used in automobiles, smart homes and so forth. Accordingly, with the increase in demand for people to converse with machines, technology that synthesizes natural speech like human speech is being actively studied.

Recently, with the development of neural networks, speech synthesis technology has made a rapid progress. Most neural speech synthesis models use a two-stage pipeline: 1) predicting a low resolution intermediate representation such as mel-spectrograms~\citep{shen2018natural, ping2017deep, li2019neural} or linguistic features~\citep{oord2016wavenet} from text, and 2) synthesizing raw waveform audio from the intermediate representation~\citep{oord2016wavenet, oord2018parallel, prenger2019waveglow, Kumar2018melgan}. The first stage is to model low-level representations of human speech from text, whereas the second stage model synthesizes raw waveforms with up to 24,000 samples per second and up to 16 bit fidelity. In this work, we focus on designing a second stage model that efficiently synthesizes high fidelity waveforms from mel-spectrograms.

Various work have been conducted to improve the audio synthesis quality and efficiency of second stage models. WaveNet~\citep{oord2016wavenet} is an autoregressive (AR) convolutional neural network that demonstrates the ability of neural network based methods to surpass conventional methods in quality. However, because of the AR structure, WaveNet generates one sample at each forward operation; it is prohibitively slow in synthesizing high temporal resolution audio. Flow-based generative models are proposed to address this problem. Because of their ability to model raw waveforms by transforming noise sequences of the same size in parallel, flow-based generative models fully utilize modern parallel computing processors to speed up sampling. Parallel WaveNet~\citep{oord2018parallel} is an inverse autoregressive flow (IAF) that is trained to minimize its Kullback-Leibler divergence from a pre-trained WaveNet called a teacher. Compared to the teacher model, it improves the synthesis speed to 1,000 times or more, without quality degradation. WaveGlow~\citep{prenger2019waveglow} eliminates the need for distilling a teacher model, and simplifies the learning process through maximum likelihood estimation by employing efficient bijective flows based on Glow~\citep{kingma2018glow}. It also produces high-quality audio compared to WaveNet. However, it requires many parameters for its deep architecture with over 90 layers.

Generative adversarial networks (GANs)~\citep{goodfellow2014generative}, which are one of the most dominant deep generative models, have also been applied to speech synthesis. \citet{Kumar2018melgan} proposed a multi-scale architecture for discriminators operating on multiple scales of raw waveforms. With sophisticated architectural consideration, the MelGAN generator is compact enough to enable real-time synthesis on CPU. \citet{yamamoto2020parallel} proposed multi-resolution STFT loss function to improve and stabilize GAN training and achieved better parameter efficiency and less training time than an IAF model, ClariNet~\citep{ping2018clarinet}. Instead of mel-spectrograms, GAN-TTS~\citep{binkowski2019high} successfully generates high quality raw audio waveforms from linguistic features through multiple discriminators operating on different window sizes. The model also shows fewer FLOPs compared to Parallel WaveNet. Despite the advantages, there is still a gap in sample quality between the GAN models and AR or flow-based models.

We propose HiFi-GAN, which achieves both higher computational efficiency and sample quality than AR or flow-based models. As speech audio consists of sinusoidal signals with various periods, modeling the periodic patterns matters to generate realistic speech audio. Therefore, we propose a discriminator which consists of small sub-discriminators, each of which obtains only a specific periodic parts of raw waveforms. This architecture is the very ground of our model successfully synthesizing realistic speech audio. As we extract different parts of audio for the discriminator, 
we also design a module that places multiple residual blocks each of which observes patterns of various lengths in parallel, and apply it to the generator.

HiFi-GAN achieves a higher MOS score than the best publicly available models, WaveNet and WaveGlow. It synthesizes human-quality speech audio at speed of 3.7 MHz on a single V100 GPU. We further show the generality of HiFi-GAN to the mel-spectrogram inversion of unseen speakers and end-to-end speech synthesis. Finally, the tiny footprint version of HiFi-GAN requires only 0.92M parameters while outperforming the best publicly available models and the fastest version of HiFi-GAN samples 13.44 times faster than real-time on CPU and 1,186 times faster than real-time on single V100 GPU with comparable quality to an autoregressive counterpart. 

Our audio samples are available on the demo  web-site\footnote{\href{https://jik876.github.io/hifi-gan-demo/}{https://jik876.github.io/hifi-gan-demo/}}, and we provide the implementation as open source for reproducibility and future work.\footnote{\href{https://github.com/jik876/hifi-gan}{https://github.com/jik876/hifi-gan}}
